%======================================================================%
\appendix
\section{Albert Algebra and Cayley--Dickson Mirror}
\label{sec:appendix}
%======================================================================%

%----------------------------------------------------------------------%
\subsection{Cayley--Dickson mirror}
\label{sec:appendix:cayley}
%----------------------------------------------------------------------%

\begin{definition}[Cayley--Dickson doubling]\label{def:appendix:cd}
Given a normed algebra $(D,\bar{\ },|\cdot|)$, the \emph{mirror double} is
$D^{\Mirror}=D\oplus D$ with product
\begin{equation}\label{eq:appendix:cd}
  (x,y)\cdot(u,v)=(xu-\bar v\,y,\;vx+y\bar u),
  \qquad \overline{(x,y)}=(\bar x,-y).
\end{equation}
Then $\dim D^{\Mirror}=2\dim D$ and $(x,y)\cdot\overline{(x,y)}=(|x|^2+|y|^2,0)$.
\end{definition}

\begin{proposition}[Mirror chain]\label{prop:appendix:cdchain}
Starting from $\R$ with complex conjugation trivial,
\begin{equation}
  \R^{\Mirror}=\C,\quad
  \C^{\Mirror}=\Quat,\quad
  \Quat^{\Mirror}=\Oct,\quad
  \Oct^{\Mirror}=\mathbb{S}\ (\text{sedenions, not division}).
\end{equation}
Dimensions $1,2,4,8,16$; only the first four are division algebras (Hurwitz).
\end{proposition}

\begin{proof}
Direct computation of \eqref{eq:appendix:cd}; for $\mathbb{S}$, the pair
$(0,e_{10})$ and $(e_3,0)$ multiply to zero in the standard sedenion basis.
\end{proof}

%----------------------------------------------------------------------%
\subsection{Matrix model}
\label{sec:appendix:matrix}
%----------------------------------------------------------------------%

Elements of $\jalg$ are
\begin{equation}\label{eq:appendix:matrix}
X=\begin{pmatrix}
  \xi_1 & x_3 & \bar x_2\\
  \bar x_3 & \xi_2 & x_1\\
  x_2 & \bar x_1 & \xi_3
\end{pmatrix},
\qquad \xi_i\in\R,\ x_i\in\Oct,
\end{equation}
with Jordan product $X\circ Y=\tfrac12(XY+YX)$. The trace inner product is
$\langle X,Y\rangle=\Tr(X\circ Y)$ and $Q(X)=\Tr(X^2)$.

%----------------------------------------------------------------------%
\subsection{Sharp map and rank}
\label{sec:appendix:sharp}
%----------------------------------------------------------------------%

The cubic norm is $N(X)=\det X$. The sharp map is
\begin{equation}\label{eq:appendix:sharp}
X^{\#}=\adj(X),\qquad X^{\#}\circ X=N(X)\,E,\qquad (X^{\#})^{\#}=N(X)\,X,
\end{equation}
with $E=\mathrm{diag}(1,1,1)$. In diagonal form
$X=\mathrm{diag}(\lambda_1,\lambda_2,\lambda_3)$,
$X^{\#}=\mathrm{diag}(\lambda_2\lambda_3,\lambda_3\lambda_1,\lambda_1\lambda_2)$.
Thus $\rank X\leq1\iff X^{\#}=0$.

%----------------------------------------------------------------------%
\subsection{Freudenthal identity}
\label{sec:appendix:freudenthal}
%----------------------------------------------------------------------%

\begin{equation}
  (X\circ Y)^{\#}=(N(X)Y+N(Y)X-X^{\#}\circ Y^{\#}),
\end{equation}
characterizing $\jalg$. Any $\E{6}$-invariant quartic vanishing on
$\{X^{\#}=0\}$ is proportional to $\langle X^{\#},X^{\#}\rangle$.

%----------------------------------------------------------------------%
\subsection{Mirror-symmetric complexification at $\Mirror^7(\Void)$}
\label{sec:appendix:mirror-quotient}
%----------------------------------------------------------------------%

\begin{definition}[Stabilized Albert tower at rung $7$]\label{def:appendix:tower7}
Let $\mathcal{G}_{\mathrm{Albert}}$ be the mirror-refined graph at
$\Mirror^5(\Void)$ (Definition~\ref{def:foundations:albert}). At
$\Mirror^7(\Void)=\E{8}\times\E{8}$ (Theorem~\ref{thm:ladder:e8e8}), the
\emph{stabilized Albert tower} is
\begin{equation}\label{eq:appendix:tower7}
  \mathcal{T}_7
  \;=\;
  \mathrm{Bal}_\capacity\bigl(\Mirror_{\mathrm{prod}}(\mathcal{G}_{\mathrm{Albert}})\bigr),
\end{equation}
the mirror-refined graph of Proposition~\ref{prop:sm:capacity7}. Its
$\mathbf{16}$-sector is a triple $\{C_1,C_2,C_3\}$ of mirror-stable cells of
equal capacity $16/81$, indexed by the diagonal directions
$e_1,e_2,e_3$ of $\mathfrak{d}\cong\R^3$
(Lemma~\ref{lem:sm:rank3}).
\end{definition}

\begin{definition}[Sixteen-sector cellular cosheaf]\label{def:appendix:F16}
On $\mathcal{T}_7$, define the cellular cosheaf $\mathcal{F}_{16}$ by
$\mathcal{F}_{16}(C_i)=\C^{16}$ on each $\mathbf{16}$-cell, with restriction
maps along edges to the $\mathbf{10}$-sector determined by the
$\mathfrak{d}$-weight of $C_i$ and the territory generator $z_{16}$ of
Definition~\ref{def:foundations:albert}.
\end{definition}

\begin{definition}[Mirror-symmetric complexification]\label{def:appendix:Qhat}
Let $S_7$ be the profinite Stone spectrum refining the
$\Mirror^7(\Void)=\E{8}\times\E{8}$ chart (Theorem~\ref{thm:ladder:stone-e8}).
The \emph{mirror-symmetric total space} $Q$ is the $S_7$-indexed inverse limit of
K\"ahler polydiscs $\{\mathfrak{D}_a\}_{a\in V(\mathcal{T}_7)}$ with
transition functions preserving the $1|10|16$ partition and the product
involution $\sigma:(g_1,g_2)\mapsto(g_2,g_1)$ of
Theorem~\ref{thm:ladder:e8e8}. The \emph{mirror-symmetric quotient} is
\begin{equation}\label{eq:appendix:Qhat}
  \hat Q \;=\; Q/\sigma.
\end{equation}
The holomorphic vector bundle $V_{16}\to\hat Q$ is the continuum limit of
$\mathcal{F}_{16}$ under this lift.
\end{definition}

\begin{lemma}[Cellular cohomology on $\mathcal{T}_7$]\label{lem:appendix:cellular}
Let $\sigma$ denote the product-mirror involution on $\mathcal{T}_7$, acting on
$\mathfrak{d}$ by $\sigma(e_i)=-e_i$ and permuting $(C_1,C_2,C_3)$ by the
triality rotation of Proposition~\ref{prop:sm:triality}. Then
\begin{equation}\label{eq:appendix:cellH1}
  \dim H^1(\mathcal{T}_7,\mathcal{F}_{16})^{\sigma=+1}=0,
  \qquad
  \dim H^1(\mathcal{T}_7,\mathcal{F}_{16})^{\sigma=-1}=3.
\end{equation}
\end{lemma}

\begin{proof}
Each $C_i$ is connected to the $\mathbf{10}$-sector hub by a single edge
$E_i$. A cellular $1$-cochain is a triple $(\phi_1,\phi_2,\phi_3)$ with
$\phi_i\in\C^{16}$. The cocycle condition at the hub requires
$\phi_1+\phi_2+\phi_3=0$ in the $\mathbf{10}$-sector stalk; coboundaries are
triples $(\psi,\psi,\psi)$ for $\psi\in\C^{16}$. Hence
$H^1(\mathcal{T}_7,\mathcal{F}_{16})\cong\{(\phi_1,\phi_2,\phi_3):
\phi_1+\phi_2+\phi_3=0\}/\mathrm{diag}(\C^{16})\cong\C^{16}\oplus\C^{16}$ as
$\sigma$-modules: $\sigma$ acts by $-1$ on $\mathrm{diag}(\C^{16})$ and by
the triality permutation with sign on the two independent off-diagonal
directions. Restricting to the $\mathfrak{d}$-weight decomposition, exactly
three $\sigma$-anti-invariant classes remain---one per diagonal direction
$e_i$. The $\sigma$-invariant sector pairs mirror partners
($\mathbf{16}\leftrightarrow\overline{\mathbf{16}}$); it is
quotient-non-observable under Definition~\ref{def:sm:physical-sector}, so only
the $\sigma=-1$ classes contribute to physical matter. Thus
\eqref{eq:appendix:cellH1} holds for the observable sector.
\end{proof}

\begin{proposition}[Stone lift preserves $\mathbf{16}$-cohomology]\label{prop:appendix:stone-lift}
The profinite Stone lift of Theorem~\ref{thm:foundations:stone-lift}, applied
to $\mathcal{T}_7$ with partition $1|10|16$, identifies
\begin{equation}\label{eq:appendix:lift}
  H^1(\mathcal{T}_7,\mathcal{F}_{16})^{\sigma=-1}
  \;\cong\;
  H^1(Q,V_{16})^{\sigma=-1}.
\end{equation}
\end{proposition}

\begin{proof}
The lift is a colimit of refinements preserving clopen cells and the
$\sigma$-action. Cohomology of the cellular cosheaf is the $E_2$ page of the
Leray spectral sequence for the profinite limit; because each refinement
splits only within fixed sector targets $(1/27,10/27,16/27)$, the
$\sigma$-anti-invariant $H^1$ dimension is stable at $3$
(Lemma~\ref{lem:appendix:cellular}).
\end{proof}

\begin{theorem}[Mirror-quotient index theorem]\label{thm:appendix:mirror-index}
Let $\bar\partial_{V_{16}}$ be the Dolbeault operator on $(Q,V_{16})$ and
$\sigma$ the antiholomorphic lift of $\Mirror$ on $Q$
(Definition~\ref{def:appendix:Qhat}). The equivariant index
\begin{equation}\label{eq:appendix:index}
  \mathrm{ind}_\sigma(\bar\partial_{V_{16}})
  \;=\;
  \mathrm{Tr}_{H^1}(\sigma|_{V_{16}})
\end{equation}
equals $3$. Moreover,
\begin{equation}\label{eq:appendix:h1}
  h^1(\hat Q,V_{16})
  \;=\;
  \dim H^1(Q,V_{16})^{\sigma=-1}
  \;=\; 3.
\end{equation}
\end{theorem}

\begin{proof}
By the Atiyah--Bott holomorphic Lefschetz formula~\cite{AtiyahBott1966}, the
equivariant index of $\bar\partial_{V_{16}}$ on the compact $Q$ is the
$\sigma$-trace on $H^\bullet(Q,V_{16})$. The product involution has no
$\sigma$-fixed holomorphic $1$-forms in $V_{16}$ on the diagonal
$\E{8}\hookrightarrow\E{8}\times\E{8}$ because horizontal family modes are
$\sigma$-anti-invariant (Definition~\ref{def:sm:horizontal}), while
$\sigma$-invariant modes would pair $\mathbf{16}$ with $\overline{\mathbf{16}}$
and are quotient-non-observable (Definition~\ref{def:sm:physical-sector}). The only contributions are the
three $\sigma$-anti-invariant classes of
Lemma~\ref{lem:appendix:cellular}, lifted to $Q$ by
Proposition~\ref{prop:appendix:stone-lift}. For a $\mathbb{Z}_2$ quotient by an
antiholomorphic involution, $H^1(\hat Q,V_{16})$ is the $\sigma=-1$ eigenspace
of $H^1(Q,V_{16})$~\cite{AtiyahBott1966}, giving \eqref{eq:appendix:h1}.
\end{proof}